%! TEX root = ../main.tex
\documentclass[main]{subfiles}

% \renewcommand{\appendixname}{付録} %%ここを付録と変更することで「付録 A」などと変更されます。デフォルトでは「appendix A」。
% \renewcommand{\appendixtocname}{付録} %%目次のappendix開始ページに表示される文字列を変更できます。
% \renewcommand{\appendixpagename}{付録} %%\appendixpage で表示されるページにでる文字列を変更できます。
% \appendixpage %%付録開始時に \appendixpagename と書かれたページを1枚追加します。
% \addappheadtotoc %%目次に付録が始まるページを追加してくれます。

\begin{document}
\chapter{付録}
\section{概要}
ソースコードは主に訓練から評価まで行うtrain\_pretrained.pyと、モデルのあらゆる定義を行うsed\_trainer\_pretrained.pyから成る。
付録では、CMTとMixStyleの主要な実装を示す。なお、SEBBsは\url{https://github.com/merlresearch/sebbs}の実装に変更を加えず、呼び出す形で使用しているため、省略する。
その他、ソースコードは\url{https://github.com/custard-1855/dcase_base}にて公開している。\par

また、使用したライブラリは以下の通り。
\begin{itemize}
\item 深層学習フレームワーク: PyTorch\cite{Ansel_PyTorch_2_Faster_2024}
\item 音声前処理: torchaudio\cite{hwang2023torchaudio}
\item モデル実装: PyTorch Lightning\cite{Falcon_PyTorch_Lightning_2019}
...
\end{itemize}

% コードの添付
\subfile{../code/01_cmt.tex}